% !PRINTED 723
\weberpage{723}

If we denote by $\alpha$, $\beta$ the indices of $t$, then $\chi_2(t) = (-1)^\alpha\Theta^\beta;$

\weq{2}{\overset{t}{\Sigma}\,\chi_2(t) = \Sigma\,(-1)^\alpha\Theta^\beta = \varphi(\Theta)}

is an \gesperrt{integer of the field} $\Omega_\nu$, and we obtain:
\[
X_2 = \frac{i\pi}{m}\,(-1,\Theta^{-1},r)\,\varphi(\Theta).
\]

Now one has to put for $\Theta$ all the roots of the equation
\weq{3}{\Theta^{\nu'} + 1 = 0}
and obtains for $A$, with regard to \S.\,197, (14):
\weq{4}{A = 2^{1-\nu'}\overset{\Theta}{\prod}\,\varphi(\Theta).}

The sum by which the number $\varphi(\Theta)$ is defined in (2)
contains $\nu'$ terms of the form $(-1)^\alpha\Theta^\beta$, in which $\alpha$, $\beta$ are to be
determined so that
\weq{5}{(-1)^\alpha\,5^\beta \equiv t \pmod{m}}
and $t$ lies between 0 and $\nu$. Now it is at once evident that
among the exponents $\beta$ no two congruent with respect to the modulus $\nu'$ can occur. For if $\beta \equiv \beta' \pmod{\nu'}$, then
$5^\beta \equiv 5^{\beta'} \pmod{\mu}$, and from
\[
(-1)^\alpha\,5^\beta \equiv t, \quad (-1)^{\alpha'}\,5^{\beta'} \equiv t' \pmod{m}
\]
there follows:
\[
(-1)^\alpha\,t - (-1)^{\alpha'}\,t' \equiv 0 \pmod{\mu}.
\]

But since $t$ and $t'$ are absolutely smaller than $\nu = \tfrac{1}{2}\mu$, this
is only possible if $t = t'$, $\alpha = \alpha'$. Accordingly
$\beta$ runs in (2) through a complete system of residues with respect to the modulus $\nu'$.

But since $\beta$ is to be taken with respect to the modulus $\nu$, the
$\beta$ determined by (5) will turn out partly smaller, partly greater than $\nu'$; and so, if we let $\beta_1$ run through the sequence of numbers $0, 1, \dots, \nu' - 1$,
we obtain in (2) two kinds of
terms:
\weq{6}{\begin{aligned}
1.\ &(-1)^\alpha\Theta^\beta = (-1)^\alpha\Theta^{\beta_1}, & \beta = \beta_1 < \nu', \\
2.\ &(-1)^\nu\Theta^\beta = -(-1)^\alpha\Theta^{\beta_1}, & \beta = \beta_1 + \nu' \geqq \nu'.
\end{aligned}}
% [sic] in row 2 the page prints the exponent \nu where \alpha is required; reproduced as printed

Now by (5) the absolutely least residue of $(-1)^\alpha\,5^\beta$
with respect to the modulus $m$ is positive, and if we therefore determine a number
$c_\beta = \pm 1$ so that the absolutely least residue of $c_\beta\,5^\beta$
becomes positive, then in the case (6), 1. we have
\[
c_{\beta_1} = (-1)^\alpha.
\]
